\documentclass[11pt]{article}
\usepackage[T1]{fontenc}
\usepackage[utf8]{inputenc}
\usepackage{lmodern,amsmath,amssymb,amsthm,mathtools,mathrsfs}
\usepackage[margin=25mm]{geometry}
\usepackage[hidelinks,hypertexnames=false]{hyperref}
\setlength{\emergencystretch}{1em}
\newcommand{\C}{\mathbb C}\newcommand{\R}{\mathbb R}
\newcommand{\PP}{\mathcal P}
\newcommand{\norm}[1]{\left\lVert#1\right\rVert}
\newcommand{\ip}[2]{\left\langle#1,#2\right\rangle}
\DeclareMathOperator{\cond}{cond}
\title{The two original low-cutoff source energies: exact minima and finite arithmetic concentration}
\author{Split-Zero mathematical continuation}
\date{19 September 2026}
\begin{document}\maketitle

This proof evaluates equations (1)--(21) of the supplied low-cutoff
continuation \cite{lce}. The comparison used on a restricted integral is
the actual pointwise upper envelope CAI17, not merely its consequence
for whole polynomial norms. The denominator is bounded by the
whole-polynomial lower comparison followed by a restricted integral in
the Gamma reference measure. Every integral in the original arithmetic
source, its mass, and every primary multiplicity are retained.

\section{Original source, complete roots and attained lifts}

Retain the original stipulated quartet domain of CAI1 and EC1--6
\cite{cai,ec}: fixed $0<\delta<1/2$, $\gamma>2$, and complete common
order $m_0\ge1$, with
\begin{equation}\tag{LCS1}
\begin{gathered}
 \rho=\tfrac12+\delta+i\gamma,\qquad
 h(s)=\prod_{\zeta\in\{\rho,\bar\rho,1-\rho,1-\bar\rho\}}
                      (s-\zeta)^{m_0},\qquad v_h=2\xi/h,\\
 w_h(y)=|v_h(1/2+iy)|^2/(2\pi),\qquad
 dm_k(y)=m_k(y)\,dy=w_h^{*k}(y)\,dy,\qquad
 \mu_h=\int_\R w_h(y)\,dy,\\
 k\equiv1\pmod4,\quad k\ge9,\quad c=k/2,\quad
 e=1+k(m_0-1),\quad q=e(k+1)^2,\\
 \beta_{ab}=c+(2a-k)\delta+i(2b-k)\gamma,\qquad
 \chi_k(S)=\prod_{a,b=0}^k(S-\beta_{ab})^e,\\
 \lambda=k\rho=c+d+it,\quad d=k\delta>0,\quad t=k\gamma>0,\quad
 z=(\lambda-c)/i=t-id,\quad R=|z|=k\sqrt{\delta^2+\gamma^2}.
\end{gathered}
\end{equation}
The total source mass is $\mu_h^k$ by Tonelli. Evenness follows from
$\xi(1-s)=\xi(s)$ and $h(1-s)=h(s)$ on the original line, and persists
under convolution. The coordinate is always $S=c+iy$. Inner products
are conjugate-linear in the first entry.

Put
\begin{equation}\tag{LCS2}
\begin{gathered}
 Q(y)=i^{-q}\chi_k(c+iy)
   =\prod_{a,b=0}^k
      [y-(2b-k)\gamma+i(2a-k)\delta]^e,\\
 \chi_{k,\lambda}(S)=\chi_k(S)/(S-\lambda)^e,\qquad
 b_\lambda=\chi_{k,\lambda}(\lambda)^{-1},\\
 L_\lambda(S)=b_\lambda\chi_k(S)/(S-\lambda),\qquad
 v_\lambda=[L_\lambda]\in E=\C[S]/(\chi_k).
\end{gathered}
\end{equation}
The $\beta_{ab}$ are distinct before their displayed repetitions, so
$b_\lambda\ne0$. Multiplication gives
$(S-\lambda)L_\lambda=b_\lambda\chi_k$. At the $\lambda$ primary,
$L_\lambda$ has leading nonzero terminal term
$(S-\lambda)^{e-1}$, and at every other primary it vanishes to order
$e$. Consequently $v_\lambda\ne0$ and $Av_\lambda=\lambda v_\lambda$.
This retains the power $e-1$ when $e>1$; no derivative of $\chi_k$ at
a repeated root is used.

The complete roots of $Q$ are closed under conjugation by
$(a,b)\mapsto(k-a,b)$ and under negation by
$(a,b)\mapsto(k-a,k-b)$. Hence $Q$ is real and even. Every root has
modulus at most $R$, all multiplicities are $e$, and none is real
because $k$ is odd and $\delta>0$. In particular, pairing all roots
with their opposites gives the pointwise inequalities
\begin{equation}\tag{LCS3}
 |Q(y)|^2\le(y^2+R^2)^q\quad(y\in\R),\qquad
 |Q(y)|^2\ge(y^2-R^2)^q\quad(|y|>R).
\end{equation}
Indeed each of the $q/2$ pairs contributes $y^2-\zeta^2$;
$y^2-R^2\le|y^2-\zeta^2|\le y^2+R^2$ on $|y|>R$, and the upper
inequality holds everywhere. Squaring the product proves LCS3.

Let $p_n(S)=i^nQ_n((S-c)/i)$ be the original monic orthogonal
polynomials, with $Q_n$ real and
$Q_n(-y)=(-1)^nQ_n(y)$, and set
$\omega_n=\norm{p_n}_{m_k}^2>0$, $b_n=[p_n]$.
For $N\ge q-1$ the remainder map and its minimum section are
\begin{equation}\tag{LCS4}
\begin{gathered}
 J_N:\PP_N\longrightarrow E,\quad J_NP=[P],\qquad
 \ker J_N=\chi_k\PP_{N-q},\\
 \mathscr S_N=H_N^{-1}J_N^*G_N,\quad
 G_N=(J_NH_N^{-1}J_N^*)^{-1}.
\end{gathered}
\end{equation}
Here $\PP_j=\C[S]_{\le j}$, $\PP_{-1}=0$, and $H_N$ is the
Gram of $1,S,\ldots,S^N$ in the original source. The remainder map
is onto, so its displayed Gram is positive definite. Direct
multiplication proves $J_N\mathscr S_N=1$ and
$\mathscr S_N(E)\perp\ker J_N$. For any other lift its difference
from $\mathscr S_N$ lies in that kernel, and the two squared norms
add. This proves both the minimum and the exact isometry from
$(E,G_N)$ to the attained source subspace.

Define the literal integrals
\begin{equation}\tag{LCS5}
\begin{gathered}
 \nu_0=\int_\R Q(y)^2m_k(y)\,dy,\qquad
 \nu_1=\int_\R y^2Q(y)^2m_k(y)\,dy,\qquad
 d\mu_Q(y)=\nu_0^{-1}Q(y)^2m_k(y)\,dy,\\
 a_z=\int_\R\frac{d\mu_Q(y)}{y-z},\qquad
 m_z=\int_\R\frac{d\mu_Q(y)}{|y-z|^2},\qquad
 s_2=\int_\R y^2d\mu_Q(y)=\nu_1/\nu_0,\\
 d_0=\omega_q/\nu_0,\qquad d_1=\omega_{q+1}/\nu_1,\qquad
 W_N=[b_N,b_{N+1},v_\lambda],\qquad
 F_N=W_N^*G_NW_N,\quad E_N=(F_N)_{33}.
\end{gathered}
\end{equation}
The auxiliary probability in LCS5 expresses ratios of original
integrals; $\nu_0$, $\nu_1$, $\omega_n$ and $E_N$ are their original
full source values throughout.

The monic minimum property gives $0<d_0,d_1\le1$. In fact both
inequalities are strict. For a positive real density the monic
orthogonal polynomial of degree $n$ has $n$ simple real roots:
otherwise, multiply its real sign-changing factors to obtain a
real polynomial of degree below $n$ whose product with the
orthogonal polynomial has constant nonzero sign on the real line,
contradicting orthogonality. Equality $d_0=1$ would make $Q$ that
orthogonal polynomial, although all its roots are nonreal.
Equality $d_1=1$ would do the same for $yQ$, which still has those
nonreal roots. The positive density needed here follows from the
original envelope displayed in LCS13 below.

\subsection{The two exact low-cutoff minima}

At $N=q-1$ there is no relation in the admitted source. Polynomial
division, monicity, parity and orthogonality give
\begin{equation}\tag{LCS6}
\begin{gathered}
 \mathscr S_{q-1}v_\lambda(c+iy)
        =b_\lambda i^{q-1}\frac{Q(y)}{y-z},\qquad
 E_{q-1}=|b_\lambda|^2\nu_0m_z,\\
 \mathscr S_{q-1}b_{q-1}=p_{q-1},\qquad
 \mathscr S_{q-1}b_q=p_q-\chi_k,\\
 (F_{q-1})_{11}=\omega_{q-1},\qquad
 (F_{q-1})_{13}=b_\lambda\omega_{q-1},\\
 (F_{q-1})_{22}=\nu_0-\omega_q,\qquad
 (F_{q-1})_{23}=i b_\lambda\nu_0a_z.
\end{gathered}
\end{equation}
For example $\ip{\chi_k}{L_\lambda}=-i b_\lambda\nu_0a_z$,
whereas $\ip{p_q}{L_\lambda}=0$ since $\deg L_\lambda=q-1$;
this proves the displayed complex phase of $(F_{q-1})_{23}$.
The coefficient of $p_{q-1}$ in $L_\lambda$ is its leading
coefficient $b_\lambda$, proving $(F_{q-1})_{13}$.

At $N=q$ the complete relation space is $\C\chi_k$. Projecting
onto precisely this line gives
\begin{equation}\tag{LCS7}
\begin{gathered}
 \mathscr S_qv_\lambda(c+iy)
   =b_\lambda i^{q-1}Q(y)\left((y-z)^{-1}-a_z\right),\\
 E_q=|b_\lambda|^2\nu_0(m_z-|a_z|^2),\qquad
 \mathscr S_qb_q=p_q-d_0\chi_k,\\
 (F_q)_{11}=\omega_q(1-d_0),\qquad
 (F_q)_{13}=i b_\lambda\omega_qa_z.
\end{gathered}
\end{equation}
Here $\ip{\chi_k}{p_q}=\omega_q$ and
$\ip{\chi_k}{L_\lambda}/\nu_0=-i b_\lambda a_z$.
Since $q$ is even, $p_{q+1}$ and $\chi_k(S)(S-c)$ have the same
odd parity in $S-c$. Their difference has degree at most $q-1$
and is orthogonal to $\chi_k$. It is consequently the attained
lift at $N=q$ of $b_{q+1}$. In particular the monic relation
$\chi_k(S)(S-c)$ has no missing preceding constant-relation
projection. This yields
\begin{equation}\tag{LCS8}
\begin{gathered}
 \mathscr S_qb_{q+1}=p_{q+1}-\chi_k(S)(S-c),\qquad
 (F_q)_{22}=\nu_1-\omega_{q+1}=\nu_1(1-d_1),\\
 \int_\R\frac{y}{y-z}\,d\mu_Q(y)=1+za_z,\qquad
 (F_q)_{23}=b_\lambda\nu_0(1+za_z).
\end{gathered}
\end{equation}
The last phase follows from
$\overline{i^{q+1}}i^{q-1}=-1$ and the preceding minus sign in
the attained lift. The term involving $a_z\int y\,d\mu_Q$ is zero.

The exact original energies are therefore
\begin{equation}\tag{LCS9}
\boxed{\begin{aligned}
 p_{1,q-1}&=\frac{\omega_{q-1}}{\nu_0m_z},&
 p_{2,q-1}&=\frac{|a_z|^2}{(1-d_0)m_z},\\
 p_{1,q}&=\frac{d_0|a_z|^2}{(1-d_0)(m_z-|a_z|^2)},&
 p_{2,q}&=\frac{|1+za_z|^2}
 {s_2(1-d_1)(m_z-|a_z|^2)}.
\end{aligned}}
\end{equation}
They equal $|(F_N)_{j3}|^2/[E_N(F_N)_{jj}]$ with exactly the
original $G_N$. The complete-primary factor cancels only between
these matching numerator and denominator terms. The variance
$m_z-|a_z|^2$ is positive because $(y-z)^{-1}$ is not constant
on a positive-density real interval. Also
\begin{equation}\tag{LCS10}
\begin{gathered}
 \operatorname{Im}a_z=-dm_z<0,\qquad
 1+za_z=\int_\R\frac{y^2}{y^2-z^2}\,d\mu_Q(y),\\
 \operatorname{Im}(1+za_z)
       =-2td\int_\R\frac{y^2}{|y^2-z^2|^2}\,d\mu_Q(y)<0.
\end{gathered}
\end{equation}
Thus every displayed response denominator and both response
numerators are nonzero.

\section{The original arithmetic envelopes and finite norm intervals}

The constants in the following display are exactly CAI16--18 and
EC14--15 \cite{cai,ec}, with different subscripts for the Gamma
density constants to avoid confusing them with unrelated scalar
coefficients:
\begin{equation}\tag{LCS11}
\begin{gathered}
 \alpha=\pi/2,\quad p_h=8m_0-\tfrac92,\quad B_h=42+8m_0,
 \quad M_h=21+4m_0,\\
 \vartheta_h=\int_{-1}^1w_h(y)\,dy>0,\quad
 M_\alpha=\int_\R e^{\alpha y}w_h(y)\,dy\in(0,\infty),\\
 c_G=\sqrt2e^{-7/6},\quad C_G=\sqrt{2\sqrt5}e^{2/3},\quad
 \sigma(y)=|\Gamma(1/4+iy/2)|^2/(2\pi),\\
 \underline a_k=
 \frac{c_h\vartheta_h^{k-3}e^{-\alpha(k-3)}(k-2)^{-B_h}}
      {C_G2^{B_h/2}},\qquad
 \overline a_k=\frac{C_h^{\rm tilt}}{c_G}
       k^{p_h+1}M_\alpha^{k-1},\\
 D_n=[1+4(n+M_h)^2]^{M_h},\qquad
 \ell_k=\underline a_k/D_{2q},\quad
 u_k=\overline a_k,\quad \mathscr L_k=\log(u_k/\ell_k).
\end{gathered}
\end{equation}
The positive $c_h$ and $C_h^{\rm tilt}$ are the original TW7 and
BT12 constants retained in CAI17. Their proved defining envelopes
and the Gamma envelopes are
\begin{equation}\tag{LCS12}
\begin{aligned}
 m_k(y)&\ge c_h\vartheta_h^{k-3}
       e^{-\alpha(|y|+k-3)}(k-2+|y|)^{-B_h},\\
 m_k(y)&\le kC_h^{\rm tilt}M_\alpha^{k-1}
       e^{-\alpha|y|}(1+|y|/k)^{-p_h},\\
 c_Ge^{-\alpha|y|}(1+|y|)^{-1/2}
       &\le\sigma(y)\le C_Ge^{-\alpha|y|}(1+|y|)^{-1/2}.
\end{aligned}
\end{equation}
These are established properties of the stipulated arithmetic
source, not new premises selected to establish the present result.
Since $p_h>1/2$ and $1+|y|/k\ge(1+|y|)/k$, the upper bounds imply
\begin{equation}\tag{LCS13}
 0<\underline a_k(1+y^2)^{-M_h}\sigma(y)
       \le m_k(y)\le u_k\sigma(y)\qquad(y\in\R).
\end{equation}
For the lower inequality use
$k-2+|y|\le(k-2)(1+|y|)$ and
$(1+|y|)^{B_h}\le2^{B_h/2}(1+y^2)^{M_h}$ in LCS12;
the extra factor $(1+|y|)^{-1/2}$ in the Gamma upper envelope
only decreases the proposed lower expression.

The Gamma reference family is Meixner--Pollaczek. Its formulas
are those in DLMF Chapter~18, authored by T.~H.~Koornwinder,
R.~Wong, R.~Koekoek, R.~F.~Swarttouw and W.~P.~Reinhardt
\cite{dlmf18}. Equations 18.19.8--9, at the parameters
$\lambda_{\rm MP}=1/4$, $\varphi=\pi/2$ and $x=y/2$, give the
exact transport
\[
 Q_n^\sigma(y)=n!P_n^{(1/4)}(y/2;\pi/2),\qquad
 \sigma(y)\,dy=\pi^{-1}w^{(1/4)}(x;\pi/2)\,dx.
\]
In particular $Q_n^\sigma$ is monic and the transported mass is
$\sqrt{2\pi}$. Their generating function, equation 18.23.7, becomes
\[
 \sum_{n\ge0}Q_n^\sigma(y)\frac{u^n}{n!}
 =(1-iu)^{-1/4+iy/2}(1+iu)^{-1/4-iy/2}
 =(1+u^2)^{-1/4}e^{y\arctan u},\quad |u|<1.
\]
Differentiation gives $(1+u^2)\partial_u\mathcal G=(y-u/2)\mathcal G$;
comparison of coefficients proves
$Q_{n+1}^\sigma=yQ_n^\sigma-n(n-1/2)Q_{n-1}^\sigma$, also the
specialization of DLMF 18.22.8. DLMF's notes attribute the generating
function to M.~E.~H.~Ismail (2009), equation 5.9.3, the recurrence
to 5.9.1, and the weight and norms to 5.9.8--9. These are
DLMF-reported book ancestry; the underlying book pages were not
newly inspected here.

For clarity, the corresponding whole-polynomial lower comparison
also follows directly. The transported orthonormal recurrence is
\[
 ye_j=\sqrt{(j+1)(j+1/2)}e_{j+1}
           +\sqrt{j(j-1/2)}e_{j-1}.
\]
Each of the raising and lowering parts has norm at most $N+1$
on polynomials of degree at most $N$, so
$\norm{yP}_\sigma\le2(N+1)\norm P_\sigma$.
Iteration and binomial expansion show
$\int|P|^2(1+y^2)^{M_h}\sigma\le D_N\norm P_\sigma^2$.
Cauchy--Schwarz with the two weights $(1+y^2)^{\pm M_h/2}$,
then LCS13, gives
\begin{equation}\tag{LCS14}
 \ell_k\norm P_\sigma^2\le\norm P_{m_k}^2
       \le u_k\norm P_\sigma^2\quad(\deg P\le2q),\qquad
 \mathscr L_k=O_h(k+\log(q+1)).
\end{equation}
The last estimate follows by taking logarithms in LCS11; all
$h$-dependent exponents and constants are fixed and $q$ is a
polynomial in $k$ of degree two or three.

For any complex monic polynomial $P$ of degree $n$ with all roots
of modulus at most $R$, define
\begin{equation}\tag{LCS15}
\begin{aligned}
 \mathscr L_n(R)&=
 \frac{2c_Ge^{-\alpha R}}{\sqrt{R+2}}
       \alpha^{-2n-1/2}\Gamma(2n+\tfrac12,\alpha),\\
 \mathscr U_n(R)&=2C_Ge^{\alpha R}
       \alpha^{-2n-1}\Gamma(2n+1).
\end{aligned}
\end{equation}
Then
\begin{equation}\tag{LCS16}
 \mathscr L_n(R)\le\norm P_\sigma^2\le\mathscr U_n(R).
\end{equation}
For the upper bound use $|P(y)|\le(|y|+R)^n$, discard the
factor $(1+|y|)^{-1/2}\le1$, put $s=|y|+R$, and extend its
integration domain to $s\ge0$. For the lower bound integrate
both real half-lines $|y|\ge R+1$, use
$|P(y)|\ge(|y|-R)^n$, put $s=|y|-R\ge1$, and use
$1+R+s\le(R+2)s$. These operations give exactly LCS15.

The Gamma monic norm in these original mass conventions is
\begin{equation}\tag{LCS17}
 \gamma_n=\sqrt{2\pi}\,n!(1/2)_n
         =\sqrt{2\pi}\,\Gamma(2n+1)/4^n.
\end{equation}
Indeed the monic version of the displayed recurrence has coefficient
$n(n-1/2)$; pairing it against consecutive monic polynomials gives
$\gamma_n=n(n-1/2)\gamma_{n-1}$, starting at the original Gamma
mass $\gamma_0=\sqrt{2\pi}$.
The monic minimum and LCS14 imply
$\ell_k\gamma_n\le\omega_n\le u_k\gamma_n$.
Consequently LCS16, applied to the full polynomials
$Q/(y-z)$, $Q$, and $yQ$, proves all three finite intervals
\begin{equation}\tag{LCS18}
\begin{gathered}
 \frac{\ell_k}{u_k}\frac{\gamma_n}{\mathscr U_n(R)}
 \le x_n\le
 \frac{u_k}{\ell_k}\frac{\gamma_n}{\mathscr L_n(R)},\\
 (n,x_n)=(q-1,p_{1,q-1}),\ (q,d_0),\ (q+1,d_1).
\end{gathered}
\end{equation}
No root is removed except for the single factor $y-z$ explicitly
present in the literal terminal eigenclass.

Here is an elementary control of the logarithmic errors in LCS18.
The classical Euler integral and Gamma recurrence are attributed
to R.~A.~Askey and R.~Roy's DLMF Chapter~5, equations 5.2.1 and
5.5.1 \cite{dlmf5}; the incomplete integrals are explicitly defined
below. For $n\ge1$, shifting the incomplete Gamma integral gives
$e^{-\alpha}\Gamma(2n+1/2)\le\Gamma(2n+1/2,\alpha)
\le\Gamma(2n+1/2)$. Further,
\[
 \frac{\sqrt\pi}{4n+1}\le
 \frac{\Gamma(2n+1/2)}{\Gamma(2n+1)}
 =\sqrt\pi\,\frac{\binom{4n}{2n}}{4^{2n}}\le\sqrt\pi.
\]
The inequalities follow because the central binomial coefficient
is the largest of the $4n+1$ coefficients whose sum is $4^{2n}$.
Substitution in LCS15--18, retaining
$a_0=\log(4/\pi)>0$, proves
\begin{equation}\tag{LCS19}
\begin{aligned}
 \log p_{1,q-1}&=-2(q-1)a_0+O_h(k+\log q),\\
 \log d_0&=-2qa_0+O_h(k+\log q),\\
 \log d_1&=-2(q+1)a_0+O_h(k+\log q).
\end{aligned}
\end{equation}
In particular $d_0,d_1$ vanish exponentially. This proof controls
the incomplete Gamma term itself and requires no unquantified
asymptotic replacement of its lower endpoint.

\section{Finite concentration in the original arithmetic measure}

Put $Y=Y_q=2q/\alpha=4q/\pi$. On the finite guard $Y>R$, LCS3
and the \emph{complete} lower comparison give
\begin{equation}\tag{LCS20}
\begin{aligned}
 \nu_0&\ge\ell_k\int_\R Q(y)^2\sigma(y)\,dy\\
 &\ge\ell_k\int_{Y\le|y|\le Y+1}Q(y)^2\sigma(y)\,dy
 \ge D_*:=
 \frac{2\ell_kc_G(Y^2-R^2)^q e^{-\alpha(Y+1)}}{\sqrt{Y+2}}.
\end{aligned}
\end{equation}
Thus $D_*$ does not assert
$\int_AQ^2m_k\ge\ell_k\int_AQ^2\sigma$ for a restricted set
$A$. The restricted integral in the second line is entirely in
the reference measure.

For $0<b<1$ write
$A_b=\{y:(1-b)Y\le|y|\le(1+b)Y\}$, and for $j=0,2$ define
\begin{equation}\tag{LCS21}
\begin{aligned}
 U_j(b)=2u_kC_G\sum_{r=0}^q\binom qr R^{2(q-r)}
 \alpha^{-(2r+j+1/2)}
 \bigl[&\gamma(2r+j+\tfrac12,\alpha(1-b)Y)\\
       +&\Gamma(2r+j+\tfrac12,\alpha(1+b)Y)\bigr].
\end{aligned}
\end{equation}
Here $\gamma(s,x)=\int_0^x u^{s-1}e^{-u}\,du$ and
$\Gamma(s,x)=\int_x^\infty u^{s-1}e^{-u}\,du$.
The pointwise upper inequality in LCS13 is valid on each part of
$A_b^c$. Apply it there, use LCS3, expand $(y^2+R^2)^q$, and
use $(1+y)^{-1/2}\le y^{-1/2}$ for $y>0$. Every coefficient is
nonnegative, all endpoint integrals are finite, and direct
integration proves
\begin{equation}\tag{LCS22}
 \int_{A_b^c}|y|^jQ(y)^2m_k(y)\,dy\le U_j(b),\qquad
 \int_{A_b^c}|y/Y|^j\,d\mu_Q(y)\le\frac{U_j(b)}{Y^jD_*}.
\end{equation}
This proves the literal finite enclosure (13) of \cite{lce} from
an actual density inequality.

\subsection{A proved exponential bound for the finite tails}

To evaluate the tails without an unsupported Laplace assertion,
define $n_j=2q+j$, $a_j=n_j+1$, and
\begin{equation}\tag{LCS23}
\begin{aligned}
 \widetilde U_j(b)=2u_kC_Ge^{\alpha R}\alpha^{-a_j}
 \bigl[&\gamma(a_j,\alpha((1-b)Y+R))\\
       +&\Gamma(a_j,\alpha((1+b)Y+R))\bigr].
\end{aligned}
\end{equation}
This is another upper bound for the left side of LCS22. Indeed,
for $y\ge0$,
$y^j(y^2+R^2)^q(1+y)^{-1/2}\le(y+R)^{2q+j}$.
The substitution $s=y+R$ maps the low interval into
$[0,(1-b)Y+R]$ after a positive extension, and maps the high
interval onto $[(1+b)Y+R,\infty)$. This gives LCS23 exactly.

For $a>0$ the density $u^{a-1}e^{-u}/\Gamma(a)$ has exponential
moment $(1-\tau)^{-a}$ for $\tau<1$, by the substitution
$(1-\tau)u=v$ in its defining integral. Markov's inequality with
$\tau=1-r^{-1}$, which is negative for $0<r<1$ and positive for
$r>1$, therefore proves
\begin{equation}\tag{LCS24}
 \frac{\gamma(a,ar)}{\Gamma(a)}\le e^{-a g(r)}\ (0<r<1),\qquad
 \frac{\Gamma(a,ar)}{\Gamma(a)}\le e^{-a g(r)}\ (r>1),\quad
 g(r)=r-1-\log r.
\end{equation}
For the lower-tail case the event $u\le ar$ implies
$e^{\tau u}\ge e^{\tau ar}$ because $\tau<0$, which justifies
the same Markov argument and its sign. The function $g$ is
positive away from $r=1$, as follows from $g'(r)=1-r^{-1}$.

Set
\begin{equation}\tag{LCS25}
\begin{gathered}
 r_{j,-}=\frac{2q(1-b)+\alpha R}{2q+j+1},\qquad
 r_{j,+}=\frac{2q(1+b)+\alpha R}{2q+j+1},\\
 C_{q,j}=\frac{2u_kC_Ge^{\alpha R}\alpha^{-a_j}\Gamma(a_j)}
                {Y^jD_*},\qquad
 \mathcal E_j(b)=C_{q,j}
   \left(e^{-a_jg(r_{j,-})}+e^{-a_jg(r_{j,+})}\right).
\end{gathered}
\end{equation}
On the explicit finite guards
$Y>R$ and $0<r_{j,-}<1<r_{j,+}$, LCS23--24 give
\begin{equation}\tag{LCS26}
 \int_{A_b^c}|y/Y|^j\,d\mu_Q(y)
       \le\frac{\widetilde U_j(b)}{Y^jD_*}
       \le\mathcal E_j(b).
\end{equation}
Every prefactor here can be evaluated from the original constants.
Its size has the following explicit bound:
\begin{equation}\tag{LCS27}
\begin{aligned}
 \log C_{q,j}\le{}&\mathscr L_k+\log(C_G/c_G)+\alpha R+\alpha
       +\tfrac12\log(Y+2)-\log\alpha\\
 &+1+\log(2q+j)+\frac{j^2}{2q}
       -q\log(1-R^2/Y^2).
\end{aligned}
\end{equation}
To verify it, substitution of $D_*$ gives the exact expression
\[
 C_{q,j}=\frac{u_kC_G}{\ell_kc_G}
 \frac{e^{\alpha R+2q+\alpha}\sqrt{Y+2}}{\alpha}
 \frac{(2q+j)!}{(2q)^{2q+j}}(1-R^2/Y^2)^{-q}.
\]
For an integer $n\ge1$, monotonicity of $\log x$ on the intervals
$[r,r+1]$ implies
$\log(n!)\le n\log n-n+1+\log n$.
Taking $n=2q+j$ and using
$\log(1+j/(2q))\le j/(2q)$ proves LCS27.

Now $R=O_h(k)$, $q\ge(k+1)^2$ and $R/Y=O_h(k/q)\to0$.
On the eventual guard $R^2/Y^2\le1/2$,
$-q\log(1-R^2/Y^2)\le2qR^2/Y^2=O_h(k^2/q)$.
Thus $\log C_{q,j}=O_h(k+\log q)=o_h(q)$.
For every fixed $b\in(0,1)$,
$r_{j,-}\to1-b$ and $r_{j,+}\to1+b$, so
$\min\{g(r_{j,-}),g(r_{j,+})\}$ has a strictly positive limit.
LCS26 consequently proves, for each such fixed $b$,
\begin{equation}\tag{LCS28}
 \int_{A_b^c}d\mu_Q\longrightarrow0,\qquad
 \int_{A_b^c}(y/Y)^2d\mu_Q\longrightarrow0
 \quad\hbox{exponentially in }q.
\end{equation}
The literal sum in LCS21 has the same exponential conclusion.
Indeed its integral form is
\[
 2u_kC_G\int_{[0,(1-b)Y]\cup[(1+b)Y,\infty)}
 y^{j-1/2}(y^2+R^2)^qe^{-\alpha y}\,dy.
\]
On $y\ge1$ its integrand is bounded by the shifted integrand used
in LCS23; on $0<y<1$ bound $(y^2+R^2)^q$ by $(1+R^2)^q$ and
integrate $y^{j-1/2}$. Thus the additional finite bound is
\begin{equation}\tag{LCS28a}
 \frac{U_j(b)}{Y^jD_*}\le\mathcal E_j(b)
 +\frac{u_kC_G}{\ell_kc_G}
  \frac{e^{2q+\alpha}\sqrt{Y+2}}{(j+1/2)Y^j}
  \left(\frac{1+R^2}{Y^2-R^2}\right)^q.
\end{equation}
The logarithm of the last term is at most
$O_h(k+\log q)+2q+q\log((1+R^2)/(Y^2-R^2))$.
Its logarithm divided by $q$ tends to $-\infty$ because
$(1+R^2)/(Y^2-R^2)=O_h(k^2/q^2)\to0$.
Consequently both $U_0(b)/D_*$ and $U_2(b)/(Y^2D_*)$ themselves
vanish exponentially, as asserted in the supplied continuation.
This conclusion uses the complete root polynomial and the actual
arithmetic upper envelope. LCS26 also supplies an independent finite
bound if the size of the sum in LCS21 makes it inconvenient.

\section{The three scalar integrals and evaluated energies}

One may take the finite bounds
\[
 \eta_b=\min\{1,U_0(b)/D_*,\mathcal E_0(b)\},\qquad
 \tau_b=\min\{U_2(b)/(Y^2D_*),\mathcal E_2(b)\}
\]
when the LCS25 guards hold; omitting $\mathcal E_j$ defines valid
bounds without those extra guards. Evenness gives
$a_z=z\int (y^2-z^2)^{-1}d\mu_Q(y)$.
For every real $y$,
$|y-z|\ge d$ and $|y^2-z^2|\ge2td$.
On the annulus $A_b$, triangle inequalities now give, whenever
$(1-b)Y>R$,
\begin{equation}\tag{LCS29}
\begin{gathered}
 \frac{1-\eta_b}{((1+b)Y+R)^2}\le m_z
 \le\frac1{((1-b)Y-R)^2}+\frac{\eta_b}{d^2},\\
 \left|\frac{Y^2a_z}{z}-1\right|
 \le B_b:=
 \frac{2b+b^2+R^2/Y^2}{(1-b)^2-R^2/Y^2}
       +\eta_b\left(1+\frac{Y^2}{2td}\right),\\
 (1-b)^2(1-\eta_b)\le\frac{s_2}{Y^2}
       \le(1+b)^2+\tau_b.
\end{gathered}
\end{equation}
For the middle estimate use
\[
 \left|\frac{Y^2}{y^2-z^2}-1\right|
 \le\frac{2b+b^2+R^2/Y^2}{(1-b)^2-R^2/Y^2}
 \quad(y\in A_b),
\]
and bound the same integrand by $1+Y^2/(2td)$ on $A_b^c$.
For the last estimate retain the second-moment contribution
$\tau_b$ on $A_b^c$ rather than discarding the tail.

The quantities $Y^2/d^2$ and $Y^2/(td)$ grow at most
polynomially in $k$, since $d=k\delta$, $t=k\gamma$, and $q$ has
degree two or three in $k$. For fixed $b$, their products with
$\eta_b$ therefore vanish by LCS28. First take $k\to\infty$ in
LCS29 and then let $b\downarrow0$. This proves all three limits
\begin{equation}\tag{LCS30}
 \boxed{Y^2m_z\longrightarrow1,\qquad
        Y^2a_z/z\longrightarrow1,\qquad s_2/Y^2\longrightarrow1.}
\end{equation}
The exponentially small central mass has been controlled even
after multiplication by the inverse denominators.

For explicit finite energy intervals, write $r_0=R^2/Y^2$ and
\begin{equation}\tag{LCS31}
\begin{gathered}
 M_- =\frac{1-\eta_b}{(1+b+R/Y)^2},\quad
 M_+ =\frac1{(1-b-R/Y)^2}+\frac{Y^2\eta_b}{d^2},\\
 A_-=(1-B_b)_+,\quad A_+=1+B_b,\quad
 S_-=(1-b)^2(1-\eta_b),\quad S_+=(1+b)^2+\tau_b,\\
 H_-=M_--r_0A_+^2,\qquad H_+=M_+-r_0A_-^2.
\end{gathered}
\end{equation}
On the displayed finite guards and $\eta_b<1$, $H_->0$,
LCS9 and LCS29 give
\begin{equation}\tag{LCS32}
\begin{aligned}
 \frac{r_0A_-^2}{(1-d_0)M_+}
 &\le p_{2,q-1}\le\frac{r_0A_+^2}{(1-d_0)M_-},\\
 \frac{d_0r_0A_-^2}{(1-d_0)H_+}
 &\le p_{1,q}\le\frac{d_0r_0A_+^2}{(1-d_0)H_-},\\
 \frac{(1-r_0A_+)_+^2}{(1-d_1)S_+H_+}
 &\le p_{2,q}\le\frac{(1+r_0A_+)^2}{(1-d_1)S_-H_-}.
\end{aligned}
\end{equation}
The guards hold for every sufficiently large $k$ at a fixed
sufficiently small $b>0$, as follows from LCS28 and $r_0\to0$.
For the last line use $|za_z|\le r_0A_+$.
The inputs $d_0,d_1$ themselves have the explicit finite intervals
LCS18. Thus the finite enclosures retain actual source integrals
and also admit enclosure solely by the original comparison constants.

Substituting LCS30 and LCS19 into the exact LCS9 gives
\begin{equation}\tag{LCS33}
\boxed{\begin{aligned}
 p_{1,N}&=\exp[-2q\log(4/\pi)+O_h(k+\log q)],\quad N=q-1,q,\\
 p_{2,q-1}&=\frac{\pi^2(\delta^2+\gamma^2)}{16}
                \frac{k^2}{q^2}(1+o_h(1)),\\
 p_{2,q}&=1-o_h(1),\\
 p_{1,q}&=d_0\frac{\pi^2(\delta^2+\gamma^2)}{16}
                \frac{k^2}{q^2}(1+o_h(1)).
\end{aligned}}
\end{equation}
To see the variance step explicitly,
$Y^2|a_z|^2=(R^2/Y^2)|Y^2a_z/z|^2\to0$, so
$Y^2(m_z-|a_z|^2)\to1$.
Also $za_z\to0$ and $s_2/Y^2\to1$, giving $p_{2,q}\to1$.
The one-sided notation in LCS33 is justified by the exact
projection inequality $p_{2,q}\le1$, proved below.
The logarithm of $R^2/Y^2$ is $O_h(\log q)$, which explains the
stated exponential rate for $p_{1,q}$.
Finally the same-class norm ratio is exactly, and then asymptotically,
\begin{equation}\tag{LCS34}
 \boxed{\frac{E_q}{E_{q-1}}=1-\frac{|a_z|^2}{m_z}
 =1-\frac{\pi^2(\delta^2+\gamma^2)}{16}
             \frac{k^2}{q^2}(1+o_h(1)).}
\end{equation}

\section{The original response matrix, spectrum and second response}

Reflection $P(S)\mapsto P(k-S)$ is an isometry in the even
arithmetic source. It preserves $\chi_k\PP_{N-q}$ and therefore
commutes with the attained minimum section LCS4. The two source
polynomials $p_N,p_{N+1}$ have opposite parities, as do their
attained lifts. Consequently $(F_N)_{12}=0$. Their normalized
quotient classes are orthonormal, and projection of
$v_\lambda/\sqrt{E_N}$ onto their span proves
$0<p_{1,N},p_{2,N}$ and $p_{1,N}+p_{2,N}\le1$.

For completeness, retain the exact JRE map \cite{jre}. Denote
the original quotient by $E_{\rm packet}=\C[S]/(\chi_k)$ to
distinguish it from the following scalar. At one
of the two cutoffs write $v=v_\lambda$, $E=E_N$,
$f=(F_{13},F_{23})^{\mathsf T}$, $B=[b_N,b_{N+1}]$, and
\begin{equation}\tag{LCS35}
\begin{gathered}
 A=B-vf^*/E:\C^2\longrightarrow(E_{\rm packet},G_N),\\
 S=A^*G_NA,\quad D=\operatorname{diag}(F_{13},F_{23}),\qquad
 T=D^{-1}SD^{-*}.
\end{gathered}
\end{equation}
Expanding $S=F_{[1,2],[1,2]}-ff^*/E$ and using $F_{12}=0$
proves the exact equality with its original source factor
\begin{equation}\tag{LCS36}
 \boxed{ET=\begin{pmatrix}p_1^{-1}-1&-1\\-1&p_2^{-1}-1\end{pmatrix}.}
\end{equation}
The eigenvalues of $T$ are those in the next display divided by
the original $E$; its mass has not been replaced.

The exact eigenvalues of $ET$ are
\begin{equation}\tag{LCS37}
 \lambda_\pm=\frac12\left[p_1^{-1}+p_2^{-1}-2
       \ \pm\sqrt{(p_1^{-1}-p_2^{-1})^2+4}\right].
\end{equation}
At $N=q-1$, put $a=p_1^{-1}-1$, $b'=p_2^{-1}-1$.
LCS33 implies $a/b'\to\infty$ and $b'\to\infty$.
For $a>b'$ the exact deviation is
\[
 \lambda_+=a+\Delta,\qquad\lambda_-=b'-\Delta,\qquad
 \Delta=\frac{2}{\sqrt{(a-b')^2+4}+(a-b')}
       \in(0,(a-b')^{-1}].
\]
It follows that $\lambda_+/a\to1$ and $\lambda_-/b'\to1$,
in particular that $ET$ is positive definite eventually. Thus
\begin{equation}\tag{LCS38}
\boxed{\begin{aligned}
 \lambda_{\max}(E_{q-1}T_{q-1})
    &=\exp[2q\log(4/\pi)+O_h(k+\log q)],\\
 \lambda_{\min}(E_{q-1}T_{q-1})
    &=\frac{16}{\pi^2(\delta^2+\gamma^2)}
                    \frac{q^2}{k^2}(1+o_h(1)),\\
 \lim_{k\to\infty}\frac1q\log\cond(E_{q-1}T_{q-1})
    &=2\log(4/\pi).
\end{aligned}}
\end{equation}

Finally let $P_K$ be any of the admitted $G_N$-orthogonal
observation-kernel projections. Retain
$K=W_N^*G_NP_KW_N$ with the three columns
$(b_N,b_{N+1},v)$, and define
$\theta=K_{33}/E\in[0,1]$, $V=K_{23}/F_{23}$.
The vector $x=P_Kv-\theta v$ satisfies
\[
 \ip v x_{G_N}=0,\qquad\norm x_{G_N}^2=E\theta(1-\theta).
\]
The residual second column
$r=b_{N+1}-v\overline{F_{23}}/E$ has norm squared
$F_{22}-|F_{23}|^2/E$. Applying Cauchy--Schwarz to
$\ip r x=K_{23}-\theta F_{23}$ gives the full finite bound
\begin{equation}\tag{LCS39}
 \boxed{|V-\theta|^2\le\theta(1-\theta)(p_2^{-1}-1).}
\end{equation}
At $N=q$, LCS33 implies
\begin{equation}\tag{LCS40}
 |V_q-\theta_q|\le
 \sqrt{\theta_q(1-\theta_q)(p_{2,q}^{-1}-1)}=o_h(1),
 \qquad 0\le\theta_q(1-\theta_q)\le\tfrac14.
\end{equation}
The convergence is uniform over those projections. It proves
$\operatorname{Im}V_q=o_h(1)$ and
$\operatorname{Re}V_q-\theta_q=o_h(1)$ without assigning a value
to $\theta_q$. For the actual joint product the exact residual is
$\overline U_q(V_q-\theta_q)$, whose modulus is
$|U_q|\,|V_q-\theta_q|$. No unproved bound on $U_q$ is used in
deducing LCS40, and that factor cannot be omitted when passing
to the product.

\begin{thebibliography}{9}
\bibitem{lce} Supplied continuation,
\emph{Low-cutoff energies, conductor and marked defect},
19 September 2026, equations (1)--(21).
Local source: \path{../inputs/LOW_CUTOFF_ENERGIES_CONDUCTOR_AND_MARKED_DEFECT.md}.
\bibitem{cai} \emph{The complete original action relative to its
arithmetic total}, CAI1--2 and CAI16--19. Local complete provider:
\path{work/mixed_tail_continuation_20260917/full_receiver_build/provider_scope/fragments/CAI_COMPLETE_PREREQUISITE.tex}.
The actual density envelopes are CAI17 and the whole-polynomial
comparison is CAI18.
\bibitem{ec} \emph{Retained eigenclass current}, full-multiplicity
class and source EC1--6; full-source constants EC14--15 and
complete-root polynomial in the following section. Local provider:
\path{work/multi_intake_20260919/arrival_02/providers/EC_COMPLETE.md}.
\bibitem{jre} \emph{The joint response ellipse in the original
three-column Gram}, JRE1--8 and JRE18--22. Local complete proof:
\path{work/multi_intake_20260919/arrival_02/JOINT_RESPONSE_ELLIPSE.tex}.
\bibitem{dlmf18} T.~H.~Koornwinder, R.~Wong, R.~Koekoek,
R.~F.~Swarttouw and W.~P.~Reinhardt,
\emph{Orthogonal Polynomials}, DLMF Chapter~18.
Weight and norms: \url{https://dlmf.nist.gov/18.19.E8},
\url{https://dlmf.nist.gov/18.19.E9}; recurrence:
\url{https://dlmf.nist.gov/18.22.E8}; generating function:
\url{https://dlmf.nist.gov/18.23.E7}.
Version 1.2.8 (15 September 2026), accessed 19 September 2026.
\bibitem{dlmf5} R.~A.~Askey and R.~Roy,
\emph{Gamma Function}, DLMF Chapter~5, Euler integral and recurrence:
\url{https://dlmf.nist.gov/5.2.E1},
\url{https://dlmf.nist.gov/5.5.E1}.
The asymptotic formulas of \S5.11 are not needed in this proof:
LCS19 uses the explicit integral and binomial inequalities displayed
there, and LCS27 uses the proved elementary factorial bound.
\end{thebibliography}
\end{document}
